\section{Complete mixed-source comparison and an unrestricted FULL-C-to-original coefficient criterion}
TEN scoped standard-axiom Lean reports verify the whole polynomial,
weighted-kernel, convolution, template, Cauchy, rank and source-payment
arguments below. Actual word/CI/counting identifications and the
unrestricted induction are the written assembly, not an end-to-end
Lean composition implementation. There is NO finite word proof premise.
General higher-capacity C and the GENERAL ORIGINAL conjecture remain OPEN.

\paragraph{Definitions and scope of the final implication.}
For an actual normalized increasing arm $X$, put $F=F_X$,
$G_k=\sum_{i=0}^kF_i$, and $D_n=(F_n-F_{n-1})^+$ for $n\ge2$
(with $D_0=D_1=0$). FULL C means
\[
 G_{k+1}D_n\ge G_kD_{n+1}\quad
 (0\le k\le n-3,\ D_nD_{n+1}>0).
\]
A final-deletion-closed FULL-C class contains every last-part deletion
of each of its arms, and every arm in it has FULL C. It need NOT be the
class of all increasing arms. For ordinary $p=q_Xq_Y$, FULL A2 means
unimodality of $2^n(p_n-p_{n-1})^+$ at EVERY rank. The original strict
core is
\[
 H=F_X+F_Y+(S_M-1)F_XF_Y,
\]
with relatives $F_X+(S_M-1)F_XF_Y$ and $F_Y+(S_M-1)F_XF_Y$, for EVERY
strict peak $M>\max(X,Y)$. The theorem below proves actual coefficient
unimodality for these polynomials under FULL C; ordinary log-concavity
of their cumulative coefficients is not substituted for that conclusion.

Keep $\rho=3/5$, $b=35/32$, $q=6125/6144<1$, $v=175/192$,
$K_0=125/22$, and $r_t=S_t(\rho)/b^{t-1}$. Frozen whole-fold norm:
$F(\rho)+K_0\le10b^M$. Frozen whole offset contraction for all
capacities at least three and the ratio envelope $r_t\le r_{\max(4,L)}$
for $t\ge L\ge3$ are available.

\subsection{A new constant-density theorem for EVERY first missing weight at least thirteen}
Theorem: for ANY actual $Y$ with $\kappa_Y\ge13$,
\[
 s_Y=\min(\lfloor M/2\rfloor,M-B),\quad m>s_Y
       \quad\Longrightarrow\quad20\,\Delta F_Y(m)\le2^m.       \tag{1}
\]
This is a CONSTANT density bound, weaker than $m\Delta F\le2^m$
when $m>20$, and is deliberately paired below with a stronger owner gap.

Such a $Y$ has LAST ELEVEN capacities at least $2,3,\ldots,12$.
Absorb the first into the unrestricted earlier prefix. The next TWO
capacities are at least three and four. Let $G$ be the prefix after
these first three terminal capacities, with skeleton degree $M_G$.
Offset contraction and the ratio envelope give
\[
                      G(\rho)\le10b^{M_G}r_4^2.       \tag{2}
\]
There remain EIGHT capacities at least $5,6,\ldots,12$, with total
skeleton degree at least $4+5+\cdots+11=60$.

Delete the LAST capacity $T$. Then
$\Delta F_Y(m)\le F_{Y_0}(m-1)$.
Expand the other SEVEN last capacities of $Y_0$:
\[
 F_{Y_0}=\sum_{\ell=0}^6\prod_{\text{last }\ell}S_t
             +\left(\prod_{\text{last }7}S_t\right)G.          \tag{3}
\]
The positive source sum at rank $m-1$ has at most the number of
compositions of $m-1$ with at most six parts. Their cut sets have size
at most five. List each cut set in increasing order and pad with zeroes
to a five-tuple; this is an injection into $m^5$ tuples. Thus its
coefficient $U_m$ satisfies $U_m\le m^5$ for $m\ge2$.
The empty-product source has zero coefficient at the ranks used here.

For the SEVEN interval kernel, stars and bars give
\[
 [z^i]\prod_{\text{last }7}[t]\le\binom{i+6}{6},\qquad
                \binom{i+6}{6}\rho^i\le3003\rho^8.     \tag{4}
\]
The last sequence has its two maxima at indices eight and nine.
Its consecutive ratio is $\rho(i+7)/(i+1)$; this gives a whole
unbounded proof, not a chosen coefficient range. The same WHOLE
positive weighted-convolution argument as the frozen two-interval
bound now yields
\[
 \rho^m\bigl(F_{Y_0}(m-1)-U_m\bigr)
                         \le3003\rho^{16}G(\rho).    \tag{5}
\]
It is ESSENTIAL to retain $U_m$; those seven source products do NOT
all vanish. Their entire contribution is paid separately.

In the BULK $M\le2m$, the eleven-part staircase gives $M\ge67$,
so $m\ge34$, and $M_G+60\le M\le2m$. Exact rational bases are
\[
 25\cdot30030\rho^{16}r_4^2\le b^{60},\qquad
              100\cdot34^5\le2^{34}.                 \tag{6}
\]
The whole induction $100m^5\le2^m$ holds for every $m\ge34$:
$(m+1)/m\le35/34$ and $(35/34)^5<2$.
Equations (2), (5), (6) and $q^m\le1$ give the long-kernel
density at most $1/25$. The ENTIRE source sum gives density at most
$m^5/2^m\le1/100$. Their sum is $1/20$, proving (1) in the bulk.
Numerically the long-kernel constant before allocating $1/25$ is
approximately $0.03899449166$; only the exact inequality (6) is used.

In the NONBULK $m>M-B$, the old largest-factor proof applies with
smaller degree thresholds. If the last capacity interval is largest,
$O=M-B\ge1+10\cdot11/2=56$, $j=m-1\ge O$, and
$\Delta F_Y(m)/2^m\le5v^j\le1/20$, from $100v^{56}\le1$.
If the initial interval is largest, its other degree is at least
$1+\cdots+11=66$, $m\ge67$, and
$\Delta F_Y(m)/2^m\le10v^m\le1/20$, from $200v^{67}\le1$.
Both are unbounded decreasing-power statements; zero and negative
slopes cause no issue. This proves (1) for all capacities, counts, weights and initial runs.

\subsection{The remaining first-missing values force a FIVE-to-one raw owner gap}
Take a positive mixed/source-prefix recovery with rising-opposite
first missing $K\in\{13,14,15\}$. At nonnegative owner reverse rank
$r$, frozen first-source and primitive facts give
\[
 D=2r+2+e,\quad h\le D,\quad \alpha+e+3\le K,\quad
       j=\alpha+h+r+2+e.                              \tag{7}
\]
If $j-2\le s_Y$, the genuine strict opposite head already pays every
proper absolute head. Otherwise
\[
 j-2>s_Y\ge S_K=\left\lfloor\frac{1+(K-2)(K-1)/2}{2}\right\rfloor,
 \qquad(S_{13},S_{14},S_{15})=(33,39,46).
\]
The finite THREE-value arithmetic of (7), without any finite owner
enumeration, implies
\[
                               e\le2r+2.              \tag{8}
\]
This is an exact unbounded linear-integer implication, checkable by
splitting ONLY $K=13,14,15$. In particular
$h+r+1\le3r+3+e\le5(r+1)$.
The frozen ordinary Hilbert growth and mixed owner phase now give,
with $Q=q_Z(r+1)$ and $R=2q_Z(r)-q_Z(r-1)$,
\[
 (r+1)Q\le(h+r+1)q_Z(r)\le5(r+1)q_Z(r),\quad
                   R\ge q_Z(r),\quad\boxed{Q\le5R}.  \tag{9}
\]

Put $N=j+1$. The frozen WHOLE dyadic source-prefix sum pays
\[
             \text{source}\ge Q\left(1-4\,2^{-N}\Delta F_Y(N)\right).
\]
Here $j-2>s_Y$ gives $N>s_Y$, so (1) implies
$4\,2^{-N}\Delta F_Y(N)\le1/5$.
By (9), the whole source pays $Q-R=C$, hence $T\le0$, contradicting
positive recovery. Thus every positive mixed recovery at $K=13,14,15$
is inside the genuine strict-head range. The negative-rank birth is
paid by the first source unless $j\le K-2$, also inside that head.
One-part owners have only this same birth exception.
Frozen absolute one-term and relative $+1$ head payments therefore
prove ENTIRE correctly conditioned MIXED proper CLOSED S at these
last three values. No arbitrary opposite capacity/count/padding cutoff
or finite word certificate is a premise.

\subsection{Complete FULL-C-to-coefficient criterion, with no first-missing bound}
The independently verified proof above closes the count induction for EVERY final-deletion-closed FULL-C class closes
without restricting first missing weights. For an old maximum at most
twelve, use the frozen arbitrary-heavy-append theorem. Otherwise:
\begin{itemize}
\item rising-opposite $K\ge16$: frozen fixed-fourteen mixed theorem;
\item $K=13,14,15$: the complete new mixed argument above;
\item $K=10,11,12$: frozen complete source-or-head certificate for ANY owner;
\item $K\le9$: an owner of first missing at least thirteen has at least
eleven non-one parts, so frozen absolute-nine/six-part and absolute-eight/
four-part nonrecovery excludes positive recovery.
\end{itemize}
The possible one-part rising opposite needs an explicit check: it has
$K=3$, so $d=K-\alpha\le2\le e+2$. The frozen FIRST-birth lower-center
reflection immediately pays every nonnegative-rank owner curvature;
its proof uses only the first source coefficient and the nondecreasing
source prefix, not a two-part count for the opposite. The negative-rank
unit birth is also present for the large owner. For $K\ge4$ an opposite
necessarily has at least two non-one parts, so the stated absolute
eight/nine inputs apply literally. Rectangular old factors are the
frozen unconditional base.
Both-forward-rising C and both-forward-nonincreasing actual resources
cover the other phases. FULL OLD A2 supplies the nonnegative head budget.
The frozen COMPLETE closed-S-to-window theorem, INCLUDING equality and
the first boundary, gives EVERY uniform window FULL A2. The actual
terminal source starts above the new first mode, so heavy append truncates
positive differences and preserves FULL A2. Proper deletion and lower-peak/
reflection joins then give ORIGINAL strict-core and BOTH-relative coefficient
unimodality at every strict peak, plus ONE arbitrary legal heavier append.

This proves the unrestricted conditional implication
\[
 \boxed{\text{ANY hereditary actual FULL-C class}
        \Longrightarrow\text{FULL A2 and ORIGINAL unamplified coefficients}.}
\]
It does NOT prove general C. In particular it does NOT solve the
GENERAL original conjecture unless general C or a replacement for its
both-rising use is established separately.

An explicit new unbounded original family follows: start with ANY
uniform hereditary FULL-C seed family (in particular the whole capacity-four
alphabet, at most two arbitrary nonbinary parts, or fixed nonbinary weight
at most 32), and append ANY FINITE NUMBER of independent successive dominant
capacities $T_i\ge2(W_0+\sum_{v<i}T_v)$, with first at least five and
empty-seed first jump handled by the frozen one-jump base. BOTH original arms
may independently have arbitrarily many such jumps, arbitrary initial runs
and binary counts, any strict peak, and ONE additional arbitrary legal
heavier append. There is no amplification and no bound of seven jumps.
